\documentclass[11pt]{article}

\usepackage{graphicx}
\usepackage{url}
\usepackage{hyperref}
\usepackage{wasysym}

\usepackage{geometry}
\geometry{letterpaper}

% See this URL for more details about the geometry Latex package:
%   ftp://ftp.tex.ac.uk/tex-archive/macros/latex/contrib/geometry/geometry.pdf
% Try what happens when you replace letterpaper in the Latex command above
% by screen, a5paper, etc.

\renewcommand{\topfraction}{1.0}
\renewcommand{\bottomfraction}{1.0}
\renewcommand{\textfraction}{0.0}
\renewcommand{\floatpagefraction}{1.0}
\renewcommand{\dbltopfraction}{1.0}


\begin{document}

{\bf \Large
\noindent Stat 5810/6810 --- Project 1 (50 Points) --- Fr 11/02/2012 \\[0.1in]
General Instructions \\[0.1in]
Due: Th 11/15/2012, 11:59pm (by e-mail)
}

\vspace{2em}

For this project, you will work with real data files that are recorded in the context of
near--infrared spectroscopy (NIRS). NIRS is an optical imaging method
for studying brain activity. In a NIRS--based experiment, 
an array of fiber--optic probes is placed on the head of a human participant
(see Figure~\ref{Hespos_NIRS} for details).

\begin{figure}[t]
\centerline{\includegraphics[width=0.85\textwidth]{Hespos_NIRS.jpg}}
\caption{\label{Hespos_NIRS}
\small
Photos of infants who were in our experiments. The photos on the top show the
cap with chin strap and elastic bandage combination that provides a snug fit between the probes
and the infant's head. The pictures on the bottom represent the placement of the probes over the
auditory and visual cortices (on the left and right respectively). The red circles indicate the
approximate placement of the source probes and the blue circles represent the detector probes.
The black squares indicate where the NIRS system would detect blood flow between each
source--detector pairing. This particular probe arrangement consists of two sets of 3 $\times$ 3 optode
arrays comprising 24 channels that each detect changes in oxy-- and deoxyhemoglobin.
[Figure and caption taken from p.~7 of
Hespos, S. J. (2010) ``What is Optical Imaging?'', Journal of Cognition and Development 11(1):3--15.]
}
\end{figure}

While the cap that is worn is fixed for each human participant in a particular study, the
exact dimensions of the head may differ considerably among participants. Therefore,
head measurements for each participant must be obtained first. These measurements
are stored in a so--called {\tt .pos} file. 

At this time, there exists no R package that allows us to read in data from these {\tt .pos} files
for further processing in R. It is our task to write functions that read in the data
from the {\tt .pos} files and provide some basic visualization of these data.

{\bf You are allowed to work in groups of up to two students. A final report in \LaTeX\
with your figures, a summary of your results, main features of your R code you
would like to highlight,
and the R code in the appendix
must be turned in as a {\tt .pdf} file. I also need the {\tt .R} file!}

Below are the specific instructions:
\begin{enumerate}
\item The {\tt .pos} files are plain ASCII file that can be opened in any text editor.
You should make yourselves familiar with the different files that are being provided.
What is identical, what is different in those files that are provided?

\item Write a function that reads in a {\tt .pos} file and creates a list  with the
content of this file. You should use regular expressions wherever possible
to extract the content from the {\tt .pos} file.
Your function should be called in the following way:
\begin{verbatim}
  Participant1 = NIRSReadPosFile(filename ="0320NLA18F.pos", processed = TRUE)
\end{verbatim}
{\tt filename} will be a required parameter. Your function should stop and report
an error if a file with this name does not exist.

{\tt processed} is an optional paramter and it can take values TRUE and FALSE. The default
is TRUE. Everything else (other than TRUE or FALSE) 
should result in a warning, but  this should be interpreted as TRUE
and the execution of the code should be continued. Thus, only FALSE is really interpreted as FALSE.
The effect of this parameter will be described in more detail later on.

Your main list with all the data should be structured as a combination of several sub--lists and matrices.
The first sub--list should be called {\tt processed} and simply contain {\tt TRUE} or {\tt FALSE},
according to the setting of this parameter.
 The second sub--list should be called {\tt FileInfo} (with elements such as
{\tt ID}, {\tt VER}, etc.). 
The third sub--list should be called {\tt User} (with elements such as 
{\tt Name}, {\tt ID}, etc.). 
The next few sub--lists should be called {\tt Probe1, $\ldots$, Probe$N$}, and {\tt MRI}. 

This should be follow by a matrix called {\tt Markers}.
The rows should be named {\tt LeftEar, RightEar, Nasion, Back,} and {\tt Top}.
The columns should be named {\tt X, Y,} and {\tt Z}.
This structure is fixed.

This should be follow by another matrix called {\tt Probes}.
The rows should be named {\tt Probe$i$ch$j$}.
The columns should be named {\tt X} to {\tt NZ}.
This structure is variable with respect to the number of rows.

Finally, there will be two more sub--lists, called {\tt Angle} and {\tt Measure}.

You cannot assume that the file format always is exactly the same as in your sample files.
Obviously, the number of probes and channels can change. Moreover, the manufacturer
of the device may decide to write additional information into this file. Surprisingly,
the serial number of the device and date and time of the recordings are not
stored in this file. Thus, it could easily happen that a future version of
this file contains entries such as {\tt SerialNumber, Date,} and {\tt Time}.
Therefore, you cannot assume that each of the main components described
above always appears in exactly the same line of the {\tt .pos} file.
Rather, you have to assume that {\tt [User]} may occur in line 8 in one
data file and in line 10 in the next {\tt .pos} file. Thus, your code has to be flexible.

 As we do not know how future additional components will be named,
we have to obtain the names for the components of our sub--lists from the data.
Then, if there will be a future entry called {\tt SerialNo} (and not
{\tt SerialNumber}), this will nevertheless get processed properly.
The following code snippet shows how you can add names to a list after
it has been created:
\begin{verbatim}
  > xx = list(c(1:4), c("a", "b"))
  > names(xx) = c("Numbers", "Letters")
  > xx$Numbers
  [1] 1 2 3 4
  > xx$Letters
  [1] "a" "b"
\end{verbatim}

\item Let us now describe the {\tt processed} parameter in more detail.
If it is FALSE, then all data will be saved exactly as stored in the {\tt .pos} file.
This means that a missing entry after the ``$=$'' will be stored as an empty string, 
the string ``~~~18y'' for Age will be saved as is, and all numbers will also 
be saved as strings. This format will allow us to check whether we are
reading the data correctly, in particular that we do not mix up any
lines and identifiers.

Otherwise, if {\tt processed} is TRUE (or anything different from FALSE),
we will further process, i.e., transform, the data according to these rules:
\begin{itemize}
\item No entry after the ``$=$'' or an arbitrary number of spaces will be 
interpreted as missing ({\tt NA}).

\item Space at the start or end should be suppressed, i.e., ``~~~Abc~~~''
should be saved as ``Abc''.

\item For the {\tt Sex} field, entries {\tt Female} and {\tt Male} should be 
translated to {\tt F} and {\tt M}, respectively. This should happen, based
on the first letter, i.e., {\tt Femal} still should be translated to {\tt F}.
The same should be done for lower case versions of the genders.
If this field is empty, the {\tt NA}--rule from above applies. If this field
contains anything else that cannot be transformed into {\tt F} or {\tt M},
a warning should be issued and {\tt NA} should be assigned.

\item For the {\tt Age} field, an entry such as ``~~~18y'' simply
should be translated into the number 18.

\item All apparent numbers should be translated into numeric. 
This includes data for the probes, the angles, and several 
variables early in the file. Note that some unusual
data entries exist, e.g., ``--0.0'', which apparently 
should be 0.
\end{itemize}

Extensively test your {\tt NIRSReadPosFile} function.
Also create copies of the data {\tt .pos} files and manipulate
those manually, e.g., by adding or deleting some variables.
You should also modify some of the content to check whether
your error--checking works properly.

\item Let's not forget that these data orignate from a human head.
We want to display the data in a 3--dimensional (3D) plot to check
that the measurements are plausible.

Install the R package {\it rgl} and read its documentation, in particular
for the 3D plot command. When we issue a command such as
\begin{verbatim}
  NIRSPlot3D(Participant1)
\end{verbatim}
the following should happen: 
\begin{itemize}
\item If the {\tt processed} list is {\tt FALSE}, we stop with an error message.
Nothing will be plotted.

\item Otherwise, all 3D {\tt X, Y, Z} coordinates  from {\tt Probe$i$ch$j$}
should be plotted in gray. Also, the 3D {\tt X, Y, Z} coordinates from the 
five markers should be plotted as well with slightly bigger symbols.
Select the colors from a 
qualitative color scheme from {\tt RColorBrewer} with three
different colors. Use one color for {\tt Nasion} and {\tt Back},
another color for the ears, and the third color for {\tt Top}.
\end{itemize}

For each of the provided {\tt .pos} files,
produce four snapshots for your report showing different 3D rotations.
Do not mix up the files and snapshots. It would be best to label
them similar to the corresponding data file, e.g., \verb|0320NLA18F_snap1.jpg|. 
Your figure caption
also should specify to which data file these snapshots belong.
\end{enumerate}


{\large \bf Good luck!!!}


\end{document}  

